\scanpage{8}
2. Structure\deleted{\unsure{s}} de carré combinatoire.

Le carré peut se définir, soit (dans le contexte des polygônes généraux) par un système
\[
(\underset{\text{sommets}}{S},\ \underset{\text{arêtes}}{A},\ \underset{\text{incidence}}{R})\tag*{(1)}
\]
satisfaisant les axiomes connus, soit (en utilisant l'\unsure{autom|autre} ``antipodisme'' $\underline{a}$ de $Q$) en termes de
\diagram{\begin{tikzcd}[ampersand replacement=\&,column sep=small,row sep=small]
\& S\arrow[d,"{\text{degré }2}"]\&\\
\text{(2)}\&\underset{(\simeq S/\underline a)}{D}\&\substack{\text{ens des diagonales}\\\text{(de card. 2)}},
\end{tikzcd}}
% slanted
\margin{Ou encore comme un ens. $S$ de cardinal $4$, avec une involution donnée sans pt fixe i.e. une \unsure{partition} de type $2,2$}
\deleted{\unsure{déf}} comme une famille de deux ensembles (les diagonales) de cardinal $2$ \deleted{\unsure{s}}, soit dualement \deleted{\unsure{pour|par}} en termes de
\diagram{\begin{tikzcd}[ampersand replacement=\&,column sep=small,row sep=small]
\& A\arrow[d,"{\text{degré }2}"]\&\\
\text{(3)}\&\underset{(\simeq A/\underline a)}{C}\&\substack{\text{ens des}\\\text{codiagonales}}
\end{tikzcd}}
% slanted
\margin{\unsure{id|ou ?}}
comme une \uline{autre} famille de deux ens. (les codiagonales \inserted{ou paires de \deleted{\unsure{côtés}} \unsure{côtés} opposés}) de cardinal $2$ \deleted{\unsure{s}}. Chacune des descriptions (1) (2) (3) est valable pour décrire le carré -- la première est \deleted{symétrique} \inserted{autoduale}, les \deleted{\unsure{deuxième|dernières|dimensions}} \inserted{deux autres} ne le sont pas, l'une privilégiant le point de vue ``sommets'', l'autre le point de vue ``arêtes'', comme ens de base (aussi dans l'un et l'autre cas c'est un ens. à quatre éléments avec involution ss \unsure{points|pts}
% archive sheet number: 7
